\section{Algorithm Details}
\label{sec:supp_algorithms}

\subsection{Dynamic Mask Thresholding via Binary Search}
\label{sec:supp_dynamic_threshold}

When an external reference mask $M_{\mathrm{ref}}$ is available (e.g., PIE-Bench ground-truth segmentation masks), we automatically determine the binarization threshold $\tau_k$ via binary search, rather than using a fixed percentile. The criterion is that the attention-derived mask should exhibit higher coverage \emph{inside} the reference region than \emph{outside}.

Concretely, let $\hat{a}_k^{\uparrow}$ denote the filtered attention map $\hat{a}_k$ resized to the spatial resolution of $M_{\mathrm{ref}}$. At each binary search step, we construct a candidate mask $M_{\mathrm{cand}} = \mathbb{1}[\hat{a}_k^{\uparrow} \geq \tau_{\mathrm{mid}}]$ and measure the inside/outside coverage ratios:
\EqCoverageRatio
If $r_{\mathrm{in}} > r_{\mathrm{out}}$, the mask is still concentrated inside the reference region — we tighten the threshold ($\tau_{\mathrm{lo}} \leftarrow \tau_{\mathrm{mid}}$). Otherwise, we relax it ($\tau_{\mathrm{hi}} \leftarrow \tau_{\mathrm{mid}}$). The search converges to the maximum threshold $\tau^*$ at which the mask still concentrates within the editing region, avoiding both over-segmentation (threshold too low, leaks into background) and under-segmentation (threshold too high, misses the target). The procedure is given in Algorithm~\ref{alg:dynamic_threshold} below.

\begin{algorithm}[t]
\caption{Adaptive Mask Thresholding via Coefficient of Variation}
\label{alg:adaptive_threshold}
\small
\begin{algorithmic}[1]
\Require Filtered attention $\hat{a}_k$, CV range $[\mathit{CV}_{\min}, \mathit{CV}_{\max}]$, multiplier range $[\kappa_{\min}, \kappa_{\max}]$
\Ensure Focus threshold $\tau_k$, preserve threshold $\tau_k^{\mathrm{low}}$

\State $\mu \gets \mathrm{mean}(\hat{a}_k)$;\quad $\sigma \gets \mathrm{std}(\hat{a}_k)$ \Comment{first two moments}
\State $\mathit{CV} \gets \sigma\,/\,(\mu + \varepsilon)$ \Comment{coefficient of variation}
\State $\kappa \gets \kappa_{\max} - (\mathit{CV} - \mathit{CV}_{\min}) \cdot (\kappa_{\max} - \kappa_{\min})\,/\,(\mathit{CV}_{\max} - \mathit{CV}_{\min})$
\State $\kappa \gets \mathrm{clamp}(\kappa,\; \kappa_{\min},\; \kappa_{\max})$ \Comment{low CV $\to$ large $\kappa$}
\State $\tau_k \gets \mu + \kappa \cdot \sigma$ \Comment{focus threshold}
\State $\tau_k^{\mathrm{low}} \gets \mu - \kappa \cdot \sigma$ \Comment{preserve threshold}
\State \Return $\tau_k,\; \tau_k^{\mathrm{low}}$
\end{algorithmic}
\end{algorithm}


\subsection{Complete Editing Pipeline}
\label{sec:supp_pipeline}

Algorithm~\ref{alg:editing} gives the full three-phase procedure corresponding to Section~4.3 in the main text.

\begin{algorithm}[t]
\caption{Three-Phase Attention-Masked Editing}
\label{alg:editing}
\small
\begin{algorithmic}[1]
\Require Source image $I_{\mathrm{src}}$, prompts $P_s, P_t$, focus word sets $\mathcal{F}_s, \mathcal{F}_t$,
         structural scales $N$, total scales $S$
\Ensure Edited image $I_{\mathrm{edit}}$

\State $\{r_k^s\}_{k=1}^{S} \gets \mathcal{Q}_{\mathrm{BSQ}}\!\bigl(\mathcal{E}(I_{\mathrm{src}})\bigr)$

\For{$k = 1$ \textbf{to} $S$}

    \Statex \textit{\quad Phase 1: Source-Stream Mask Extraction}
    \State Generate tokens at scale $k$ conditioned on $P_s$
    \If{$k \ge N$}
        \State Compute $a_{k,b}$ for focus words $\mathcal{F}_s$; filter via IQR $\to \hat{a}_k^{(s)}$
        \State $\tau_k^{(s)},\, \tau_k^{\mathrm{low},(s)} \gets \mathrm{AdaptiveThreshold}(\hat{a}_k^{(s)})$ \Comment{Alg.~\ref{alg:adaptive_threshold}}
        \State $M_{k,s}^{\mathrm{f}} \gets \mathbf{1}[\hat{a}_k^{(s)} \ge \tau_k^{(s)}]$,\quad
               $M_{k,s}^{\mathrm{p}} \gets \mathbf{1}[\hat{a}_k^{(s)} \le \tau_k^{\mathrm{low},(s)}]$
    \EndIf

    \Statex \textit{\quad Phase 2: Target-Stream Mask Extraction (with Background Anchoring)}
    \State Generate tokens at scale $k$ conditioned on $P_t$
    \If{$k \ge N$}
        \State Anchor positions where $M_{k,s}^{\mathrm{p}}{=}1$ to $r_k^s$ \Comment{hard background lock}
        \State Compute $a_{k,b}$ for focus words $\mathcal{F}_t$; filter via IQR $\to \hat{a}_k^{(t)}$
        \State $\tau_k^{(t)},\, \_ \gets \mathrm{AdaptiveThreshold}(\hat{a}_k^{(t)})$
        \State $M_{k,t}^{\mathrm{f}} \gets \mathbf{1}[\hat{a}_k^{(t)} \ge \tau_k^{(t)}]$
        \State $M_{k,t}^{\mathrm{f}} \gets M_{k,t}^{\mathrm{f}} \;\wedge\; \neg\, M_{k,s}^{\mathrm{p}}$ \Comment{preserve overrides target focus}
    \EndIf

    \Statex \textit{\quad Phase 3: Mask-Guided Token Composition}
    \State $M_k^{\mathrm{rep}} \gets \neg\!\bigl(M_{k,s}^{\mathrm{f}} \cup M_{k,t}^{\mathrm{f}}\bigr)$
    \State Generate $g_k$ at scale $k$ conditioned on $P_t$
    \If{$k < N$}
        \State $r_k \gets r_k^s$ \Comment{full structural anchoring}
    \Else
        \State $r_k \gets M_k^{\mathrm{rep}} \odot r_k^s + (1{-}M_k^{\mathrm{rep}}) \odot g_k$
    \EndIf

\EndFor

\State \Return $\mathcal{D}_{\mathrm{VAE}}\!\bigl(\textstyle\sum_{k} \mathrm{up}(\mathrm{dec}(r_k))\bigr)$
\end{algorithmic}
\end{algorithm}

